\documentclass[12pt,a4paper]{article}

\usepackage[T1]{fontenc}
\usepackage[utf8]{inputenc}
\usepackage[brazilian]{babel}
\usepackage{amsmath}
\usepackage{booktabs}
\usepackage{geometry}
\usepackage{graphicx}
\usepackage{tabularx}
\usepackage{xcolor}
\usepackage[hidelinks]{hyperref}

\geometry{margin=2.5cm}
\setlength{\parindent}{1.25cm}
\setlength{\parskip}{0.2em}

\title{DeliveryDispatch\\
\large Motor Híbrido de Busca e Monitoramento Geoespacial\\
\large Relatório do Trabalho Prático 1 --- DCC207}
% Preencher com todos os integrantes confirmados antes da entrega.
\author{}
\date{Outubro de 2026}

\begin{document}

\maketitle
\thispagestyle{empty}

\begin{abstract}
Este relatório descreve o desenvolvimento do DeliveryDispatch, uma solução para
indexar restaurantes e uma frota sintética de entregadores em Belo Horizonte. A
abordagem combina codificação Geohash, uma Trie compacta implementada no projeto,
geração de posições por distribuição normal e projeção dos pontos sobre a malha
viária do OpenStreetMap. Além das decisões algorítmicas, são apresentados o fluxo
de preparação dos dados e a estratégia de testes.
\end{abstract}

\tableofcontents
\clearpage

\section{Introdução}

O problema proposto consiste em construir um motor de busca e monitoramento
geoespacial para uma frota de entregadores em operação na cidade de Belo
Horizonte. A partir de um restaurante selecionado, o sistema deve localizar
rapidamente os entregadores candidatos, eliminar aqueles fora de um raio de
atendimento e apresentar os resultados em um mapa interativo, acompanhados de
distâncias e estimativas de tempo de chegada \cite{enunciado}.

Uma busca exata sobre toda a frota teria custo proporcional ao número total de
entregadores em cada consulta. Para reduzir esse espaço, a solução associa cada
posição a uma célula Geohash e usa prefixos dessas células como chaves de uma
Trie compacta. A seleção por prefixo fornece um conjunto candidato; em seguida,
uma etapa geométrica deve tratar células adjacentes e verificar a distância real
até o restaurante. Essa combinação busca conciliar velocidade de indexação com
precisão espacial.

O desenvolvimento ocorre no contexto do Trabalho Prático 1 da disciplina
DCC207 --- Algoritmos 2. Este documento concentra a descrição técnica e as
decisões de projeto, enquanto o README do repositório permanece voltado à
instalação e à execução.

\section{Arquitetura e fluxo de dados}

A implementação foi dividida em algoritmos, modelos, acesso a dados e serviços.
O fluxo principal pode ser resumido nas etapas a seguir:

\begin{enumerate}
    \item leitura da base fornecida de bares e restaurantes;
    \item geocodificação dos endereços ainda pendentes;
    \item persistência do CSV enriquecido e construção dos objetos de domínio;
    \item geração de $N$ posições sintéticas ao redor dos estabelecimentos;
    \item projeção dessas posições para ruas transitáveis;
    \item codificação das coordenadas em Geohash;
    \item inserção de entregadores, nomes e endereços em Tries compactas.
\end{enumerate}

A Tabela~\ref{tab:modules} relaciona as principais responsabilidades aos módulos
que as implementam.

\begin{table}[htbp]
\centering
\small
\caption{Organização dos componentes principais.}
\label{tab:modules}
\begin{tabularx}{\textwidth}{@{}p{0.41\textwidth}X@{}}
\toprule
Módulo & Responsabilidade \\
\midrule
\path{algorithms/geohash.py} &
Codificação e decodificação Geohash \\
\path{algorithms/compact_trie.py} &
Trie compacta genérica e busca por prefixo \\
\path{data/bar_loader.py} &
Leitura, validação, retomada e persistência do CSV \\
\path{services/geocoding.py} &
Cliente sequencial do Nominatim com cache local \\
\path{services/courier_generation.py} &
Distribuição espacial e construção dos entregadores \\
\path{services/road_network.py} &
Carga da malha viária e projeção dos pontos \\
\path{main.py} &
Orquestração do fluxo e construção dos índices \\
\bottomrule
\end{tabularx}
\end{table}

Os modelos \texttt{Bar} e \texttt{Courier} são dataclasses imutáveis. Essa
decisão torna explícito quando um registro é substituído após a geocodificação e
evita alterações acidentais enquanto o mesmo objeto é referenciado pelos índices.

\section{Preparação da base geográfica}

\subsection{Leitura e persistência}

O carregador valida a presença das colunas obrigatórias e converte cada linha em
um objeto \texttt{Bar}. O CSV processado acrescenta latitude, longitude e
\texttt{GEOCODING\_STATUS}; o Geohash é calculado em memória ao carregar o
registro e, portanto, não é duplicado no arquivo.

A persistência usa um arquivo temporário, posteriormente movido sobre o destino.
Esse procedimento reduz o risco de deixar um CSV parcialmente escrito em caso de
interrupção. Durante uma execução de geocodificação, também são salvos
\emph{checkpoints} a cada 100 consultas. Registros com estado
\texttt{found} ou \texttt{not\_found} não são consultados novamente, permitindo
retomar o processamento dos itens ainda marcados como \texttt{pending}.

\subsection{Geocodificação e cache}

O cliente Nominatim envia as requisições sequencialmente, identifica a aplicação
por \emph{User-Agent}, limita a frequência a uma requisição por segundo e
mantém um cache JSON local. Essas medidas reduzem consultas repetidas e seguem
as principais restrições do serviço público \cite{nominatim-policy}. O cache não
é versionado.

Ao término do processamento, a base apresentou a distribuição
da Tabela~\ref{tab:geocoding}.

\begin{table}[htbp]
\centering
\caption{Resultado da preparação da base fornecida.}
\label{tab:geocoding}
\begin{tabular}{@{}lrr@{}}
\toprule
Estado & Registros & Proporção \\
\midrule
Encontrados (\texttt{found}) & 19\,238 & 92,23\% \\
Não encontrados (\texttt{not\_found}) & 1\,621 & 7,77\% \\
Pendentes (\texttt{pending}) & 0 & 0,00\% \\
\midrule
Total & 20\,859 & 100,00\% \\
\bottomrule
\end{tabular}
\end{table}

\section{Algoritmos e decisões de projeto}

\subsection{Codificação Geohash}

A implementação inicia com os intervalos $[-180,180]$ para longitude e
$[-90,90]$ para latitude. Cada bit divide um dos intervalos ao meio. Os bits de
índice par refinam a longitude e os de índice ímpar refinam a latitude, portanto
a intercalação começa pela longitude. A cada cinco decisões binárias, o valor
obtido seleciona um caractere do alfabeto Base32 do Geohash.

Para uma precisão de $p$ caracteres, são realizadas $5p$ subdivisões. Assim, o
tempo e o espaço da codificação são $O(p)$; neste projeto, $1 \leq p \leq 8$.
A decodificação percorre o caminho inverso e retorna o centro da célula
representada. Como referência de escala, no equador uma célula com oito
caracteres mede aproximadamente $38{,}2\,\mathrm{m} \times 19\,\mathrm{m}$,
com dimensões que variam conforme a latitude \cite{elastic-geohash-grid}.

Geohash fornece uma aproximação espacial, não uma distância métrica. Além disso,
pontos próximos podem cair em células diferentes quando estão em lados opostos
de uma borda.

\subsection{Trie compacta}

A Trie foi implementada no próprio projeto, sem bibliotecas prontas para a
estrutura central. Cada nó guarda o trecho completo da aresta que o conecta ao
pai, em vez de armazenar apenas um caractere. Os filhos ficam em um dicionário
indexado pelo primeiro caractere da aresta. Uma folha especial identificada por
\texttt{\$} marca o fim da chave e armazena um valor genérico.

Na inserção, há três casos principais: criação de uma nova aresta, continuação
por uma aresta totalmente compatível ou divisão da aresta no maior prefixo
comum. A remoção executa o processo inverso e volta a unir caminhos que ficam
com um único filho não terminal. Esse mecanismo reduz cadeias de nós sem
bifurcação.

A busca consome o prefixo ao longo das arestas compactadas e, depois da
correspondência, coleta as folhas da subárvore. Considerando acesso médio
constante ao dicionário de filhos, a navegação custa $O(m)$ e a materialização
dos resultados custa o número de nós visitados. No contrato esperado, isso é
expresso como $O(m+k)$, em que $m$ é o tamanho do prefixo e $k$ representa o
volume retornado.

\subsection{Índices de nome e endereço}

Nomes e endereços de restaurantes são normalizados com conversão para caixa
minúscula, decomposição Unicode, remoção de acentos e redução de sequências de
espaços. Duas Tries separadas recebem esses textos e os respectivos objetos
\texttt{Bar}, preparando a busca tanto por nome quanto por endereço.

\section{Busca geoespacial e filtragem}

\section{Geração sintética de entregadores}

A frota é gerada somente a partir dos estabelecimentos que possuem coordenadas.
Depois de uma embaralhação reproduzível dos centros, os $N$ entregadores são
atribuídos em rodízio aos restaurantes. Dessa forma, a diferença entre as
quantidades associadas a quaisquer dois centros utilizados é de, no máximo, uma
unidade.

Para cada entregador, são sorteados dois deslocamentos independentes,
$X,Y \sim \mathcal{N}(0,\sigma^2)$, com $\sigma=500\,\mathrm{m}$. Os
deslocamentos são convertidos aproximadamente de metros para graus por

\begin{align}
\Delta\varphi &=
    \frac{X}{111\,320}, \\
\Delta\lambda &=
    \frac{Y}{111\,320\max\left(|\cos\varphi|,10^{-12}\right)},
\end{align}

em que $\varphi$ é a latitude do restaurante. A correção pelo cosseno compensa
a variação do comprimento de um grau de longitude. A semente padrão é 42, o que
permite repetir os mesmos sorteios em testes e experimentos.

A geração matemática tem custo $O(N)$, desconsiderando a projeção viária. Após
essa projeção, cada entregador recebe um identificador sequencial, as
coordenadas corrigidas, o identificador do restaurante usado como centro e seu
Geohash.

\section{Correção para a malha viária}

Pontos produzidos por uma distribuição contínua podem cair dentro de edifícios,
quadras ou corpos d'água. Para aproximar posições válidas de circulação, o
serviço viário usa o OSMnx e dados do OpenStreetMap
\cite{osmnx-docs,osm-copyright}.

Na primeira execução, é obtido o grafo simplificado da rede de automóveis de
Belo Horizonte (\texttt{network\_type=drive}). O grafo é projetado para um
sistema de coordenadas adequado a operações métricas e salvo em GraphML. Nas
execuções seguintes, esse arquivo é reutilizado.

As coordenadas sintéticas, inicialmente em EPSG:4326, são convertidas para o
mesmo sistema do grafo. O OSMnx encontra a aresta mais próxima de cada ponto.
Em seguida, a biblioteca geométrica interpola a posição sobre essa aresta e o
resultado retorna a EPSG:4326. Quando uma aresta não possui geometria explícita,
o segmento entre seus nós inicial e final é usado como aproximação.

Os dados cartográficos são disponibilizados pelo OpenStreetMap sob a licença
ODbL e requerem atribuição ao projeto e a seus colaboradores
\cite{osm-copyright}.

\section{Interface interativa}

\section{Testes e garantia de qualidade}

Os testes automatizados verificam os contratos das partes isoladas e da
orquestração. Entre os cenários cobertos estão:

\begin{itemize}
    \item valores conhecidos e entradas inválidas na codificação Geohash;
    \item inserção, divisão de arestas, busca, remoção e valores da Trie;
    \item leitura, persistência e retomada do CSV processado;
    \item consulta e reaproveitamento do cache de geocodificação sem rede real;
    \item equilíbrio e reprodutibilidade da geração de entregadores;
    \item projeção sobre uma aresta de um grafo sintético;
    \item integração do fluxo principal por meio de substitutos controlados.
\end{itemize}

Protocolos para o geocodificador e o projetor viário permitem usar implementações
falsas nos testes, mantendo a lógica de domínio independente da rede. A
integração contínua executa análise estática, verificação de formatação, testes e
construção do pacote com Python 3.12 em Linux e Windows.

\section{Avaliação experimental}

\section{Conclusão}

O DeliveryDispatch estabelece a base de dados e os principais componentes
algorítmicos da solução. A codificação Geohash e a Trie compacta foram
implementadas diretamente, enquanto bibliotecas geoespaciais foram reservadas
para geocodificação, sistemas de coordenadas e malha viária. A geração sintética
reproduzível e a projeção sobre ruas completam a infraestrutura necessária para
produzir uma frota plausível.

\begin{thebibliography}{9}

\bibitem{enunciado}
Vimieiro, Renato.
\emph{Trabalho Prático 1 --- Indexação Geoespacial e Geometria Computacional Híbrida}.
DCC207 --- Algoritmos 2, Universidade Federal de Minas Gerais, 2026.

\bibitem{elastic-geohash-grid}
Elastic.
\emph{Geohash grid aggregation: cell dimensions at the equator}.
Disponível em:
\url{https://www.elastic.co/docs/reference/aggregations/search-aggregations-bucket-geohashgrid-aggregation}.
Acesso em: 19 out. 2026.

\bibitem{osmnx-docs}
OSMnx.
\emph{OSMnx documentation}.
Disponível em: \url{https://osmnx.readthedocs.io/}.
Acesso em: 19 out. 2026.

\bibitem{nominatim-policy}
OpenStreetMap Foundation.
\emph{Nominatim Usage Policy}.
Disponível em:
\url{https://operations.osmfoundation.org/policies/nominatim/}.
Acesso em: 19 out. 2026.

\bibitem{osm-copyright}
OpenStreetMap.
\emph{Copyright and License}.
Disponível em: \url{https://www.openstreetmap.org/copyright}.
Acesso em: 19 out. 2026.

\end{thebibliography}

\end{document}
